\chapter*{Tóm tắt}
\addcontentsline{toc}{chapter}{Tóm tắt}

Tổng hợp ảnh đa phương thức y tế là quá trình kết hợp thông tin từ nhiều loại ảnh chụp khác nhau --- ví dụ CT, MRI, PET, SPECT --- thành một ảnh duy nhất giúp bác sĩ vừa quan sát được cấu trúc giải phẫu, vừa nắm bắt được thông tin chức năng của cơ quan. Đây là công cụ hỗ trợ chẩn đoán quan trọng trong điều trị các bệnh lý phức tạp như ung thư, đột quỵ và rối loạn thần kinh.

Mô hình CDDFuse \cite{zhao2023cddfuse} (CVPR 2023) là một trong những phương pháp dẫn đầu hiện tại, sử dụng kiến trúc dual-branch Transformer--CNN để phân rã đặc trưng ảnh thành thành phần \emph{Cơ sở} (Base --- tần số thấp, cấu trúc toàn cục) và \emph{Chi tiết} (Detail --- tần số cao, đặc trưng riêng từng modality). Tuy nhiên, CDDFuse áp dụng \emph{cùng một quy tắc tổng hợp} --- phép cộng đơn giản $f_F = f_V + f_I$ --- cho cả hai nhánh, không khai thác được sự khác biệt bản chất giữa Base và Detail.

Đồ án này đề xuất chiến lược \textbf{tổng hợp bất đối xứng} (Asymmetric Fusion): thay vì dùng cùng một quy tắc cho cả hai nhánh, thiết kế \emph{quy tắc riêng biệt} phù hợp với đặc tính của từng nhánh. Thực nghiệm được tổ chức theo ba giai đoạn:

\begin{enumerate}\itemsep2pt
    \item \textbf{Giai đoạn 1 --- Khảo sát quy tắc Cơ sở}: huấn luyện và đánh giá 5 quy tắc thay thế cho nhánh Base (L1-norm, Visual Saliency Map, Local Energy, Local Entropy, Weighted Average Scalar), giữ nguyên nhánh Detail theo paper gốc.
    \item \textbf{Giai đoạn 2 --- Khảo sát quy tắc Chi tiết}: huấn luyện và đánh giá 8 quy tắc thay thế cho nhánh Detail (Spatial Frequency, SML, L1-norm, MaxAbs, Local Energy, Adaptive Gating, Saliency Gate, Soft PCNN), giữ nguyên nhánh Base theo paper gốc.
    \item \textbf{Giai đoạn 3 --- Kết hợp bất đối xứng}: ghép các quy tắc tốt nhất từ hai giai đoạn, tạo mô hình \textbf{Comb-WAvg-SML} (Base $=$ Weighted Average Scalar, Detail $=$ Sum-Modified-Laplacian).
\end{enumerate}

Toàn bộ thực nghiệm được thực hiện trên tập dữ liệu Harvard Medical Image Fusion (738 cặp ảnh huấn luyện, 72 cặp ảnh kiểm thử trên 3 modality CT/PET/SPECT), đánh giá bằng 22 chỉ số chất lượng theo chuẩn VIFB và xếp hạng tổng hợp bằng phương pháp Composite Z-score.

\textbf{Kết quả chính}:
\begin{itemize}\itemsep2pt
    \item Mô hình đề xuất \textbf{Comb-WAvg-SML} vượt CDDFuse rõ rệt trên CT-MRI (z-score $+0{,}240$ so với $-0{,}483$, chênh lệch $+0{,}72$), đặc biệt ở các chỉ số chất lượng biên và tần số không gian (SF, AG, EI, QG).
    \item So sánh với APGFusion 2024 (kiến trúc phức tạp hơn: PoolFormer, ConvolutionalGLU, HFF block) cho thấy cải tiến kiến trúc không vượt rõ CDDFuse, khẳng định \emph{chiến lược fusion rule quan trọng hơn kiến trúc đơn thuần}.
\end{itemize}

\bigskip
\noindent\textbf{Từ khóa:} tổng hợp ảnh y tế đa phương thức, CDDFuse, quy tắc tổng hợp bất đối xứng, Composite Z-score, Modality-Specificity, Harvard Medical Image Fusion.
